\subsection*{15. Revenue Is Not Bookings, Billing, or Collections}

A \textbf{booking} generally describes a commercial contract or commitment, often a management measure rather than a standardized accounting term. \textbf{Billing} is the amount invoiced. \textbf{Collections} are cash received. \textbf{Recognized revenue} follows the applicable recognition rules for the promises fulfilled. In long-duration contracts, all 4 can happen in different months and for different amounts.

Suppose a college agrees to an annual applied-elective program for ₹18 lakh. It signs an agreement today, pays an initial ₹3 lakh next month, and requires several milestones. A founder who spends against ₹18 lakh as if it were collected cash will incur liabilities that the contract itself may not fund in time. A signed proposal that lacks budget authorization is weaker still.

Maintain a contract register that separates signed amount, contingencies and conditions precedent, invoice schedule, acceptance milestones, collected funds, performance completed, and remaining obligations. In a healthy institutional business, sales, teaching operations, invoicing, and finance look at the \emph{same contract} from different angles.

\subsection*{16. Depreciation and Capital Expenditure: Neither Is Optional Economically}

Suppose a learning studio buys laptops or equipment for ₹3 lakh. The purchase consumes ₹3 lakh of cash at once. If the asset qualifies for capitalization, accounting may spread its depreciable cost across useful periods rather than expense the full purchase immediately. This can make reported current profit look better than current cash flow, but it is not an error: the asset can support future output.

A venture must forecast both the \textbf{cash needed to acquire and replace assets} and the accounting effect of using them. A business that celebrates EBITDA while continually purchasing expensive equipment may consume far more cash than EBITDA suggests. Conversely, a 1-time equipment purchase does not necessarily mean the current business model loses money on each student.

Classify expenditures under the applicable accounting and tax rules; do not arbitrarily capitalize ordinary salaries, marketing costs, or routine content work merely to improve apparent profit. The managerial decision is to fund productive assets only when cash, utilization, maintenance, obsolescence, and quality justify them.

\subsection*{17. Debt Proceeds Are Financing, Not Sales}

A bank loan temporarily increases cash and liabilities. It does not make the company profitable. Loan principal repayments reduce cash and debt and are not normally treated as an income-statement expense. Interest creates borrowing cost; the appropriate recognition and cash-flow presentation depend on standards and facts. Equity investment increases cash and owners' capital but is not revenue either.

This distinction should be obvious, yet founders often unconsciously blend all cash inflows into 1 runway number and treat future loan capacity as if it were already certain. A pledged credit facility can be withdrawn or restricted by covenants. An investor indication is not cash. An unapproved personal loan is not enterprise funding.

Build a financing schedule with outstanding principal, interest rate and basis, payment dates, collateral, guarantees, covenants, maturity, fees, and refinancing assumptions. Mark what is \textbf{contracted and available} separately from what management hopes to obtain.

\subsection*{18. Cash Profits and Normalized Economic Profits}

In a founder-led venture, the founder may teach, sell, design, code, review, and perform administration without drawing a business salary. The bank ledger may look attractive because the company's most expensive input is being donated.

2 views are necessary:
\begin{enumerate}
\item \textbf{Cash survival view:} what will actually leave the company's account under existing arrangements?
\item \textbf{Normalized economics view:} what would it cost to replace founder work with competent paid labor at realistic intensity?
\end{enumerate}
A venture can pass the 1st and fail the 2nd. It can still be a worthwhile founder-owned small business, but the limits should be understood before promising a scalable institution.

Do not blindly assign the founder a market salary as though they already employ themselves under a standard contract. Managerial normalization is an economic test, not automatically a statutory journal entry. Chapter 17 introduced the distinction; financial survival requires tracking both views every month.

\subsection*{19. The Financial Statements Are Reports; Survival Requires Forecasts}

A balance sheet and income statement largely describe the consequences of what has happened. A cash-flow statement explains past movements. Financial survival depends on \textbf{what will happen before the next major decision point}.

A monthly income statement cannot warn the founder that rent is due Tuesday, payroll Thursday, and the customer will probably pay next month. A beautifully reconciled balance sheet does not itself tell the founder how long a supplier might extend credit. An annual budget may conceal a dangerous week of negative cash balances even if annual net cash is positive.

The missing instrument is a short-horizon, cash-dated operating forecast, combined with a longer-horizon scenario model. The next part develops those tools and the working-capital engine that drives them.
